\section{Advanced Method --- Fourth-Order Gauss Implicit Runge--Kutta}
% ============================================================

\subsection{Method Formulation}

As an additional advanced comparison, we consider the classical
two-stage, fourth-order Gauss--Legendre implicit Runge--Kutta method
described by Butcher~\cite{butcher2016}, referred to here as Gauss IRK4.
The coefficients used below are the standard two-stage Gauss coefficients
given in Chapter~3, Section~34.2 of Butcher's text.

This is a classical Gauss collocation method rather than the four-stage
classical explicit RK4 scheme. The method is included as an extension
to the three baseline methods studied in Sections~2--4.

Using the same column-wise vectorisation as in the main analysis, let
\[
y=\operatorname{vec}(X)\in\mathbb{R}^{9},
\qquad
f(y)=\operatorname{vec}\!\left[X(I-X^TX)\right].
\]

The two Gauss nodes are
\[
c_1=\frac12-\frac{\sqrt3}{6},
\qquad
c_2=\frac12+\frac{\sqrt3}{6},
\]
with weights
\[
b_1=b_2=\frac12,
\]
and coefficient matrix
\[
A=
\begin{pmatrix}
\frac14 & \frac14-\frac{\sqrt3}{6}\\[1mm]
\frac14+\frac{\sqrt3}{6} & \frac14
\end{pmatrix}.
\]

For one time step of size $h$, the stage values satisfy
\begin{equation}
\label{eq:gauss_stages}
Y_i
=
y_n+h\sum_{j=1}^{2}a_{ij}f(Y_j),
\qquad i=1,2,
\end{equation}
and the new solution is
\begin{equation}
\label{eq:gauss_update}
y_{n+1}
=
y_n+\frac{h}{2}\left[f(Y_1)+f(Y_2)\right].
\end{equation}

Because the two stages depend on each other, they must be solved as one
coupled nonlinear system rather than as two independent nine-dimensional
problems.

\subsection{Solution of the Implicit Stages}

Define the coupled residuals
\[
F_i(Y_1,Y_2)
=
Y_i-y_n-h\sum_{j=1}^{2}a_{ij}f(Y_j),
\qquad i=1,2.
\]

The $(i,j)$ block of the Newton Jacobian is
\[
\frac{\partial F_i}{\partial Y_j}
=
\delta_{ij}I_9-h\,a_{ij}J_f(Y_j).
\]

Therefore, each Newton iteration requires the solution of an
$18\times18$ linear system. The implementation uses the predictors
\[
Y_i^{(0)}
=
y_n+c_i h f(y_n),
\]
followed by damped Newton iteration. A trial correction is accepted only
when the full coupled stage residual remains finite and decreases.

The stopping condition is
\[
\|F\|_{\infty}
\le
10^{-12}\left(1+\|Y\|_{\infty}\right).
\]

This tolerance controls the residual of the implicit stage equations. It
should not be interpreted as a bound on the global error of the numerical
solution.

\subsection{Order and Convergence}

The classical two-stage Gauss collocation method has order
\[
p=2s=4.
\]

For exactly solved stage equations, the one-step defect is
$O(h^5)$, and the corresponding per-unit-step local truncation error is
$O(h^4)$. Since the vector field
\[
f(X)=X(I-X^TX)
\]
is polynomial, it is smooth on bounded sets relevant to the finite-time
integration considered here.

To connect the method with the global error framework used in
Section~2, suppose that $f$ is Lipschitz with constant $L$ on a bounded
neighbourhood containing the relevant stage states. For sufficiently
small $h$ satisfying
\[
h\|A\|_{\infty}L<1,
\]
the coupled stage equations have a locally unique solution in that
neighbourhood, and the resulting one-step map has a uniform local
Lipschitz bound. The discrete Gr\"onwall argument used previously then
gives global fourth-order convergence over a fixed finite time interval,
provided that the exact and numerical trajectories remain in the same
bounded neighbourhood.

In practice, the stage equations are solved only approximately. If the
stage residual is $\varepsilon(h)$ and the local inverse bound remains
valid, the induced stage error is $O(\varepsilon)$ and the accumulated
global contribution is also $O(\varepsilon)$. Therefore, a strict
asymptotic fourth-order statement as $h\to0$ requires the stage residual
to decrease appropriately with $h$. The implementation uses the fixed
threshold $10^{-12}$, so the numerical evidence below should be
interpreted as a finite-grid verification rather than as an independent
asymptotic theorem.

\subsection{Absolute Stability}

For the scalar test equation
\[
y'=\lambda y,
\qquad
z=h\lambda,
\]
the amplification factor of the two-stage Gauss method is
\begin{equation}
\label{eq:gauss_stability}
R_G(z)
=
\frac{1+\frac{z}{2}+\frac{z^2}{12}}
     {1-\frac{z}{2}+\frac{z^2}{12}}.
\end{equation}

Let
\[
N(z)=1+\frac{z}{2}+\frac{z^2}{12},
\qquad
D(z)=1-\frac{z}{2}+\frac{z^2}{12}.
\]
Then
\[
|D(z)|^2-|N(z)|^2
=
-2\operatorname{Re}(z)
\left(1+\frac{|z|^2}{12}\right).
\]

Hence, when $\operatorname{Re}(z)\le0$,
\[
|R_G(z)|\le1,
\]
so the method is A-stable. This is the classical A-stability property of
Gauss methods; see Butcher, Chapter~3, Section~35.3~\cite{butcher2016}.

However,
\[
R_G(z)\to1
\qquad
\text{as }
z\to-\infty
\]
along the negative real axis. Therefore, the Gauss method is not
L-stable. In particular, very strongly damped modes are not necessarily
suppressed as aggressively as they are under Implicit Euler.

For the initial frozen Jacobian of the present problem, the fastest
negative mode is approximately $\lambda=-26$. With $h=0.08$,
\[
z=-2.08,
\]
for which the Gauss amplification factor is approximately $0.1335$.
This frozen-spectrum calculation is only a local diagnostic and should
not be interpreted as a complete nonlinear stability proof.

\begin{figure}[htbp]
    \centering
    \includegraphics[width=0.78\textwidth]
    {advanced_irk4/stability_comparison.png}
    \caption{Comparison of the scalar stability amplification factors
    along the negative real axis. The dashed reference marks the initial
    frozen mode $z=-26\times0.08=-2.08$. The Gauss IRK4 curve illustrates
    its A-stability and its lack of L-stability.}
    \label{fig:gauss_stability_comparison}
\end{figure}

\subsection{Numerical Comparison}

The Gauss IRK4 implementation was tested on the same initial matrix and
time interval $[0,2]$ used in the main experiments. The terminal
Frobenius-norm error was again measured against the exact finite-time SVD
solution.

Using the refined grids
\[
N=64,\ 128,\ 256,
\]
the fitted convergence slope for Gauss IRK4 is approximately
\[
3.9674,
\]
which is consistent with the expected fourth-order behaviour on the
tested finite grids.

For comparison, the corresponding fitted slopes in the same comparison
script are approximately
\[
1.0101,\quad
3.9274,\quad
1.0137
\]
for Explicit Euler, classical RK4, and Implicit Euler, respectively.

\begin{figure}[htbp]
    \centering
    \includegraphics[width=0.78\textwidth]
    {advanced_irk4/convergence_comparison.png}
    \caption{Terminal Frobenius-norm error versus step size for the four
    numerical methods. The fitted Gauss IRK4 slope on the refined grids
    is approximately $3.9674$, consistent with fourth-order behaviour on
    the tested finite grids.}
    \label{fig:gauss_convergence_comparison}
\end{figure}

A matched-accuracy comparison was also carried out. Classical RK4 with
$N=75$ gives terminal error
\[
5.09443\times10^{-8}
\]
using 300 right-hand-side evaluations. Gauss IRK4 with $N=64$ gives
terminal error
\[
5.07898\times10^{-8},
\]
using 720 right-hand-side evaluations, 264 Jacobian evaluations, and
132 Newton iterations. The two terminal errors differ by approximately
$0.3\%$.

Thus, for this $3\times3$ benchmark and this accuracy level, Gauss IRK4
does not reduce the number of right-hand-side evaluations relative to
classical RK4. In addition, the evaluation counts do not include the cost
of Jacobian assembly, the $18\times18$ linear solves, or Python overhead,
so they should not be interpreted as direct wall-clock timing results.

\begin{figure}[htbp]
    \centering
    \includegraphics[width=0.78\textwidth]
    {advanced_irk4/cost_comparison.png}
    \caption{Terminal error versus counted right-hand-side evaluations.
    The highlighted RK4 and Gauss IRK4 runs have terminal errors differing
    by approximately $0.3\%$. The plotted work measure does not include
    Jacobian assembly, linear solves, or Python overhead.}
    \label{fig:gauss_cost_comparison}
\end{figure}

On the common grid $h=0.1$, the sampled Lyapunov energy decreases at each
stored time for all four methods. This is an observation from the present
experiment rather than a general proof that Gauss IRK4 preserves energy
monotonicity for arbitrary step sizes.

\begin{figure}[htbp]
    \centering
    \includegraphics[width=0.78\textwidth]
    {advanced_irk4/energy_comparison.png}
    \caption{Sampled Lyapunov energy trajectories for the four methods on
    the common grid $h=0.1$. All sampled trajectories decrease in this
    experiment; this is an empirical observation rather than a general
    monotonicity theorem.}
    \label{fig:gauss_energy_comparison}
\end{figure}

Finally, tightening the stage residual threshold from $10^{-12}$ to
$10^{-14}$ changes the terminal solution by approximately
\[
4.44\times10^{-16},\qquad
1.50\times10^{-13},\qquad
3.47\times10^{-14}
\]
for $N=64,128,256$, respectively. These changes are much smaller than
the corresponding time-discretisation errors, supporting the use of the
fixed Newton tolerance over the tested grid range.

The numerical values reported in this subsection are generated by the
advanced-method comparison scripts. The convergence slopes,
matched-accuracy errors, and work counts are recorded in
\texttt{data/advanced\_irk4/convergence\_cost.csv} and
\texttt{data/advanced\_irk4/summary.json}. The energy trajectories shown
in Figure~\ref{fig:gauss_energy_comparison} are recorded in
\texttt{data/advanced\_irk4/energy\_trajectories.csv}.

% ============================================================